% !TeX root = main.tex



\section{Experimental Setup}

This section describes the hardware platform, benchmark
configurations, and performance metrics used to evaluate the gPCD
framework.

\subsection{Hardware Platform}

All experiments were conducted on a Lenovo ThinkPad P16 Gen~2 mobile
workstation equipped with a 13th Generation Intel Core i9-13980HX
processor and an NVIDIA RTX 3500 Ada Generation professional GPU. The
GPU contains 5,120 CUDA cores and 12~GB of ECC GDDR6 memory with a
192-bit memory interface providing approximately 432~GB/s of memory
bandwidth.This configuration provides a modern mobile GPU baseline for evaluating real-time, 
large-scale collision detection systems~\cite{hinumNVIDIARTX3500}.

The gPCD framework was implemented using Vulkan, with broad-phase occupancy 
generation and rendering executed through the graphics pipeline and narrow-phase 
evaluation executed through compute shaders.

\subsection{Maximum Particle Capacity}

The maximum number of particles that can be represented in a single
simulation is constrained by the available storage-buffer capacity of
the graphics device. Let $M_{\max}$ denote the maximum addressable
storage-buffer range supported by the device and let $S_p$ denote the
memory footprint of a single particle. The corresponding upper bound on
the particle count is

\begin{equation}
N_{\max}=\left\lfloor\frac{M_{\max}}{S_p}\right\rfloor.
\end{equation}

This relationship was used to determine the largest particle counts
considered during testing. As the particle structure size increases,
the maximum achievable particle count decreases proportionally.
Conversely, reducing the memory footprint of the particle structure
increases the maximum problem size that can be represented within a
single storage buffer. Consequently, all test cases were selected to
remain below the storage-buffer capacity limit imposed by the target
hardware.

\subsection{Verification and Performance Evaluation}

Because gPCD performs exact collision detection, verification and
performance evaluations were conducted as separate stages. Each test
case was first executed in a verification mode to validate particle
counts, cell occupancy generation, and collision counts. Verification
runs were used to ensure that all particles were processed correctly
and that all collision events were detected and counted as expected.

Following successful verification, the same configurations were
executed in performance mode to measure execution time, throughput,
and time-per-particle scaling. Performance measurements were therefore
obtained only after the correctness of the collision-detection results
had been established. This procedure ensured that all reported
performance metrics correspond to validated collision-detection
outcomes.


\subsection{Benchmark Configuration}

Benchmark datasets consisted of randomly distributed particle
configurations with increasing particle count while maintaining an
approximately constant collision fraction. Particle counts ranged from
\numprint{32} to
\numprint{10000000} particles.

\subsection{Performance Metrics}

Performance evaluation was based on graphics execution time
($t_g$), compute execution time ($t_c$), total execution time
($t_t$), throughput, and time-per-particle.

Throughput was computed as the number of particles processed per
second,

\begin{equation}
T=\frac{N}{t_t},
\end{equation}

where $N$ is the number of particles and $t_t$ is the total
execution time.

Time-per-particle was computed as

\begin{equation}
t_{pp}=\frac{t_t}{N},
\end{equation}

which provides a normalized measure of execution cost independent of
particle count.

Additional benchmarks were performed to evaluate the effects of memory
locality and collision density on framework performance.
